\section{Two obstructions to finite and positive spectral closures}
\label{fgo:section}

These statements concern the actual free friend-tree generator
\[
 (\mathcal L f)(T)=\sum_{w\in T}h_T(w)[f(T+w)-f(T)],
 \qquad h_T(w)=\sum_{u\sim_Tw}\frac1{d_T(u)},
\]
and the nonleaf count $C(T)$. They do not assert irrationality of the
growth exponent. The first result strengthens the failure of a closure
in the core count alone to the failure of every finite-dimensional
constant-coefficient linear observable closure containing that count.

\begin{lemma}[A pole of increasing order on stars]
\label{fgo:star-pole}
Let $S_d$ be the star with $d$ leaves, for an integer $d\ge2$. For each
$r\ge0$, the values $(\mathcal L^r C)(S_d)$ are the restrictions of a
rational function $f_r(d)$. For $r\ge2$, this function has the Laurent
expansion
\[
 f_r(d)=(-1)^{r-1}d^{1-r}+O(d^{2-r})\qquad(d\to0).
\]
In particular its pole at zero has exact order $r-1$.
\end{lemma}
\begin{proof}
Write the generator as the sum of its off-diagonal birth operator and
its diagonal term $-|T|I$. Expand $\mathcal L^r C$ as a sum over words
of $r$ operations. An off-diagonal operation specifies its source vertex,
then its chosen neighbour, and contributes the reciprocal source degree.
A diagonal operation contributes minus the current population. For fixed
$r$, there are only finitely many history patterns after identifying
histories that differ by a permutation of the original star leaves.

Suppose such a pattern mentions $k$ distinct original leaves. Assigning
them distinct labels gives the falling-factorial multiplicity
$(d)^{\underline{k}}$. Degrees of vertices other than the original centre
are positive integers determined by the finite pattern. The centre has
degree $d+j$, where $j\ge0$ is the number of births attached to it.
Diagonal population factors are polynomials in $d$. Thus each pattern
has a rational weight in $d$, proving the first assertion.

Every pattern containing a birth mentions at least one original leaf:
the first birth either has the centre as source and an original leaf as
target, or has an original leaf as source and the centre as target.
Its falling-factorial multiplicity therefore has a factor $d$. The only
possible poles at $d=0$ come from centre-source choices made before any
birth has attached to that centre. If a word has at most $r-1$ such
choices, its pole order is at most $r-2$. A word with no births is
polynomial in $d$.

It remains to collect words having $r$ such choices. These are exactly
the histories consisting of $r$ successive centre-source births, with
no diagonal operations. The centre keeps degree $d$, and each event
chooses uniformly one of its $d$ original neighbours. A neighbour
becomes a nonleaf exactly when first selected. Consequently the combined
contribution of these histories is
\[
 1+d\left[1-\left(1-\frac1d\right)^r\right].
\]
Its coefficient of $d^{1-r}$ is $(-1)^{r+1}=(-1)^{r-1}$.
Every other word has smaller pole order, so this coefficient survives
in $f_r$.
\end{proof}

For example,
\[
 f_0(d)=f_1(d)=1,\qquad
 f_2(d)=\frac12-\frac1d,\qquad
 f_3(d)=\frac{d^3+11d^2+d+6}{6d^2(d+1)}.
\]

\begin{theorem}[No finite linear generator closure]
\label{fgo:independence}
The functions
\[
 C,\mathcal LC,\mathcal L^2C,\ldots
\]
are linearly independent over $\mathbb C$ on the class of finite free
trees with at least three vertices. Hence there is no finite-dimensional
vector space of observables containing $C$ and invariant under
$\mathcal L$.
\end{theorem}
\begin{proof}
If $\sum_{j=0}^R c_j\mathcal L^jC=0$, restrict to stars $S_d$ with
integer $d\ge2$. The resulting rational function vanishes at infinitely
many points and hence is identically zero. If the largest index with
nonzero coefficient were $R\ge2$, Lemma~\ref{fgo:star-pole} would give
an uncancelled pole of order $R-1$, a contradiction. Thus only
$c_0C+c_1\mathcal LC$ remains. Stars give $c_0+c_1=0$, whereas the
four-vertex path has $C=2$ and $\mathcal LC=1$, giving $2c_0+c_1=0$.
Both coefficients vanish. A finite invariant space containing $C$
would contain every generator iterate and contradict this independence.
\end{proof}

The theorem rules out a finite constant-matrix linear differential system
of observables containing $C$ that is valid from every finite seed.
It does not rule out population-dependent recurrences, nonlinear
reductions, special-function equations in the discount parameter, or
an accidental scalar relation for the edge-start mean alone.

\subsection{An exact obstruction to a positive exponential mixture}

Let $F_2(t)=\mathbb E_{\rm edge}C(t)$ and
$c_r=(\mathcal L^rC)(\mathrm{edge})$. The Yule generating-function
representation makes $F_2$ analytic in a neighbourhood of zero, and
the generator hierarchy gives $F_2^{(r)}(0)=c_r$. Direct finite exact
recursion yields
\[
 c_1=2,\quad c_2=-2,\quad c_3=4,\quad c_4=-\frac{19}{3},
 \quad c_5=\frac{133}{12},\quad c_6=-\frac{8189}{432},
 \quad c_7=\frac{70753}{2160}.
\]
Put $m_j=c_{j+1}/2$. These cannot be the moments of a positive measure
on the real line: their Hankel determinant is
\[
 \det(m_{i+j})_{0\le i,j\le3}=-\frac{43427}{3732480}<0.
\]
An equivalent certificate uses the integer polynomial
\[
 q(x)=468x^3+969x^2+19x-437.
\]
If $\mathfrak m(x^j)=m_j$, then exact rational arithmetic gives
\[
 \mathfrak m(q^2)=-\frac{564551}{40}<0.
\]
Therefore there is no positive real measure $\nu$, with finite
exponential moments in a neighbourhood of zero, for which
\[
 \frac{F_2'(t)}2=\int_{\mathbb R}e^{\lambda t}\,d\nu(\lambda)
\]
holds near zero. Such a representation would imply
$\mathfrak m(q^2)=\int q(\lambda)^2\,d\nu(\lambda)\ge0$.
The conclusion is specific to this proposed scalar positive spectral
representation; it does not rule out signed measures, non-self-adjoint
operators, or other observables.

\paragraph{A failed arithmetic transfer.}
The increasing star poles do not by themselves prove a denominator
lower bound for the edge-start derivatives. For example, a hub of
degree $5$ followed by repeated centre-source events suggests a
$5$-denominator contribution, but
$c_6=-8189/432$ has no factor $5$ in its denominator. Near-star histories
attain the same valuation and can cancel. More concretely, the
generator words whose first four operations are births ending at
$S_5$ have total prefix weight $2\cdot4!=48$; their remaining
two-generator contribution is $48f_2(5)=72/5$. Its factor $5$
disappears in the complete value of $c_6$. Thus the statewise pole
noncancellation theorem must not be used as an unproved noncancellation
claim after averaging over the full edge-start distribution.
